Can a video game teach strangers to coordinate without ever practicing together?
We study whether individual game experience transfers coordination skills to unrelated tasks with new partners---a capability we call \emph{cross-domain zero-shot coordination}.
Using large language model agents as human proxies, we design and evaluate three coordination-training games against untrained and irrelevant-game baselines across 31 coordination tasks spanning focal point selection, tacit role assignment, and convention formation.
Our AI-designed multi-skill game, \syncminds, which combines focal point training, perspective-taking, role differentiation, and convention formation, achieves the most balanced improvement profile: 2.9$\times$ improvement on individual focal point tasks, 1.98$\times$ on hard role assignment (50.5\%$\to$100\%), and 24\% faster convention convergence.
We find that single-skill training produces narrow gains and can even \emph{hurt} coordination on complementary tasks---focal point training causes agents to converge when they should differentiate, collapsing role assignment to 0\% on some tasks.
The multi-skill game resolves this convergence--differentiation tension by teaching agents \emph{when} to match and \emph{when} to divide.
Our results provide proof-of-concept that AI can design coordination-enhancing games with measurable cross-domain transfer, opening a new direction for scalable coordination training in ad-hoc teams.
